% Proposed insertions/replacements for QI_MLPs___ICLR_2027_Submission (18).pdf.
% This is a patch fragment, not a standalone document. Not applied to the paper.
% Locations below refer to printed manuscript line numbers, not source lines.
% All new figures and substantive exposition belong in the appendix.
% Main-text replacements still require compilation in the actual paper project
% to confirm unchanged pagination; no full main-text TeX source is present here.
% Keep existing equation/theorem labels when replacing their contents.

% MAIN-TEXT BUDGET: replacements only; no new figure or paragraph.
% EDIT 1. Page 4, Section 3.1: retain the first sentence at line 183
% (width and hidden parameters). Replace the following prose and Equation (4).
% TSVD denotes the truncated-SVD pseudoinverse, defined here in words;
% Appendix C supplies its threshold convention and expanded formula.
Given $M=sN+1$ equispaced observations, form the unscaled feature matrix
$A$ and sample vector $y$. Our empirical readout uses the truncated-SVD
pseudoinverse:
\begin{equation}
  \widehat c=\operatorname{TSVD}_{\tau_{\mathrm{emp}}}(A)y.
  \tag{4}
\end{equation}
Appendix~C specifies the cutoff $\tau_{\mathrm{emp}}$ and compares this
solve with the reference-scaled recovery analyzed in Appendix~A.

% EDIT 2. Page 5, Theorem 3.1: replace only the sentence at line 232.
The reference-scaled recovery of Appendix~A returns a network
$\widehat q_W$, with readout $c_W$, satisfying
% Keep Equations (5)--(6) and the rest of the theorem unchanged.
% Keep the main theorem's a_H=1 special case; the new experiment below
% explicitly uses the general Appendix A.1 geometry instead.

% EDIT 3. Page 4, Figure 2 caption: replace only its opening sentence.
Empirical performance of QUILL with unscaled truncated-SVD readouts.
% Retain the descriptions of panels (a)--(c).

% EDIT 4. Appendix C: insert a separate subsection after line 1896, after
% the existing precision experiment's Validation paragraph. This comparison
% does not share that experiment's all-operations-at-p-bits protocol.
\subsection{Comparison of recovery procedures}
\label{sec:recovery_comparison}
For clarity, the exact-arithmetic form of the empirical solve (4) is
\[
 \operatorname{TSVD}_{\tau_{\mathrm{emp}}}(A)y
 =\sum_{\sigma_k>\tau_{\mathrm{emp}}}
       \frac{u_k^\top y}{\sigma_k}\,v_k,
 \qquad A=U\Sigma V^\top.
\]
Here $\sigma_k$ are the singular values and $u_k,v_k$ the corresponding
left and right singular vectors; the cutoff is specified for each experiment.
This truncated solve need not minimize the unrestricted least-squares objective.

We compare the reference-scaled recovery procedure of Appendix~A with the
practical unscaled SVD solve on the six functions in
Figure~\ref{fig:recovery_comparison}. Both receive the same neurons and
$M=8N+1$ equispaced observations on $[-1,1]$. We fix $\lambda=1/4$,
$R=\lceil(11/8)\sqrt N\rceil$, $\delta=4$, and $\Delta=9/16$, using the
general geometry of Appendix~A.1. Thus the hidden width is $W=N+2R+1$.
We use $N\in\{64,128,196,256,320,384,448,512,576,640,704,768\}$.
The sampling parameters in (A.55) are $s=8$, $p=1$, and $\varrho=2$.
For non-dyadic widths we use the uniform input-error route (A.81).

The practical solve uses FP64 DGELSD with relative cutoff $2^{-52}$.
The reference-scaled solve uses FP64 DGEJSV, the analytic bound (A.51),
the cutoff condition (A.65), and the sharper numerical bounds given below.
The two solvers are compared as complete procedures; this is not an
ablation of column scaling alone. To preserve the nominal geometry,
we store the exactly representable affine parameters
$\gamma=N/8$ and $\beta_j=N/8-j/4$, so each neuron is
$\tanh(\gamma x+\beta_j)$ for $j=-R,\ldots,N+R$.

Target samples are prepared offline at higher precision and rounded to FP64.
The comparison uses correctly rounded tanh values, FP64 model arrays and
readout arithmetic, and pairwise output summation. Analytic and numerical
bound verification uses 256-bit directed MPFR arithmetic and FP64 expansion
products. This is separate from the preceding precision-controlled experiment.
All 72 function--resolution pairs pass the stated geometry and recovery
checks. Errors are measured on 16,384 uniform midpoint samples; doubling
the evaluation grid changes the reported errors by less than $0.68\%$.

At every tested $N\ge512$, both procedures attain measured relative errors
below $6\times10^{-13}$ on all six functions. At smaller widths the
theoretical cutoff can be substantially more conservative, and the unscaled
solver often attains high accuracy earlier. These results demonstrate high
measured accuracy for the theoretically specified procedure and support the
simpler solver as a practical implementation. They do not extend the recovery
guarantee to the unscaled solver.

\begin{figure}[t]
  \centering
  % Copy the generated comparison.pdf to this project-relative location.
  \includegraphics[width=\linewidth]{figures/recovery_comparison.pdf}
  \caption{Comparison of readout recovery procedures. Solid curves use
  unscaled SVD; dashed curves use reference-scaled SVD with the prescribed
  theoretical cutoff. Colors identify target functions. Each pair uses
  identical neurons and samples, with fixed parameters across resolutions.
  The horizontal variable $N$ counts interior intervals; the hidden width
  is $W=N+2\lceil(11/8)\sqrt N\rceil+1$. Both methods attain measured
  relative $L^2$ error below $6\times10^{-13}$ at every tested $N\ge512$;
  the unscaled procedure often attains high accuracy earlier. Errors are
  evaluation-grid measurements, not certified continuous-domain error bounds.}
  \label{fig:recovery_comparison}
\end{figure}

% EDIT 5. Appendix A.3.3, page 22: insert after line 1145, before A.3.4.
% A.63 remains a valid sufficient realization. This paragraph supplies the
% sharper realization used by the new experiment, not a changed cutoff.
\paragraph{Sharper numerical bounds for the recovery comparison.}
The comparison in Section~\ref{sec:recovery_comparison} uses the following
alternative sufficient realization of (A.61)--(A.62). Let
$H=\widehat U^\top\widehat U-I$ and
$G=\widehat V^\top\widehat V-I$ be the orthogonality defects of the
returned singular-vector matrices. Retain the bounds
$\|H\|_2\le o_U<1$ and $\|G\|_2\le o_V<1$ above, and write
$\sigma_{\max}=\|\widehat\Sigma\|_2$ and $V_+=\sqrt{1+o_V}$.
Define the scalar remainder bounds
\[
 t_s(e)=\frac{e^2}{8(1-e)^{3/2}},\qquad
 t_i(e)=\frac{3e^2}{8(1-e)^{5/2}}\quad(0\le e<1).
\]
With $e_F$ and $r_F$ as above, one may use
\[
 a_W=e_F+r_F+
 \left(\frac12\|H\widehat\Sigma\|_2+
       t_s(o_U)\sigma_{\max}\right)V_+
 +\frac12\|\widehat\Sigma G\|_2+t_s(o_V)\sigma_{\max}.
\]
Let $p$ be the computed pairwise projection from (A.60), and let
$q=\widehat U^\top\widetilde y$ denote the exact projection of the stored
arrays, used only for analysis. Define
\[
 d_p=\|p-q\|_2+\frac12\|Hq\|_2+t_i(o_U)\|q\|_2.
\]
Then a valid label perturbation bound is
\[
 b_W=e_y+d_p+u\|p\|_2+\mu\|\widehat\sigma\|_2,
\]
where $\widehat\sigma$ is the vector of returned singular values;
$e_y$, $u$, and $\mu$ retain their meanings above. All displayed norm
quantities are replaced by validated upper bounds in the implementation.

To verify these formulas, write
$\widehat U=U_0(I+H)^{1/2}$ and
$\widehat V=V_0(I+G)^{1/2}$ with orthonormal $U_0,V_0$.
Taylor's theorem and the spectral theorem give
\[
 \|(I+H)^{1/2}-I-H/2\|_2\le t_s(o_U),\qquad
 \|(I+H)^{-1/2}-I+H/2\|_2\le t_i(o_U),
\]
and the corresponding square-root bound for $G$.
Subtracting $U_0\widehat\Sigma V_0^\top$ from
$\widehat U\widehat\Sigma\widehat V^\top$ and grouping the left and
right polar corrections proves the stated $a_W$ bound. Since
$U_0^\top\widetilde y=(I+H)^{-1/2}q$, the inverse-square-root bound gives
$\|p-U_0^\top\widetilde y\|_2\le d_p$. The same retained-coordinate
construction of $y_0$ used above, together with its division-rounding
allowance, proves the stated $b_W$ bound. The export bound (A.64) and
cutoff condition (A.65) remain unchanged.

The validation computes enclosures of residual and Gram products by
decomposing FP64 inputs into fixed-point chunks. For chunk width $b$ and
inner dimension $k$, the implementation checks
$2b+\lceil\log_2 k\rceil\le52$ and representable exponent ranges,
so the digit products and sums are exact under the stated FP64 arithmetic.
Reconstruction uses compensated sums, with decomposition tails and
rounding allowances included. Norm bounds include these enclosure errors;
when exact normalization is unavailable, a validated Frobenius-norm bound
is used. Scalar bounds are rounded outward in MPFR. These calculations
validate the numerical bounds and do not refine the fitted coefficients.
